% §10 Conclusion — target 0.25 page

\section{Conclusion}
\label{sec:conclusion}

We presented \systemname, the first reproducible benchmark for evaluating LLM
penetration testing agents at \emph{network scale} on IoT infrastructure, and
\harnessname, a network-native harness that keeps LLM context tractable
through per-device and per-vulnerability parallelization. Our empirical study
of five state-of-the-art models across 7~scenarios and three architectural
patterns\todo{(update to five once firewall rules land)} shows that LLM agents
outperform traditional scanners on configuration and credential categories
while underperforming on industrial protocols, and that the
detection-to-exfiltration (L1--L3) evidence spectrum exposes model
differences that binary-flag benchmarks conflate.

\paragraph{Future work}
(i)~\emph{Credential propagation} — closing the Phase~4 $\to$ Phase~2 loop to
enable multi-hop exploitation chains;
(ii)~\emph{Real hardware} — extending to embedded Zigbee, LoRaWAN, and
sub-GHz attack surfaces via SDR tooling;
(iii)~\emph{Defender's side} — using \systemname{} as a benchmark for
LLM-assisted IoT defense (detection, patching recommendation).

All code, scenarios, ground truth, and run logs are released at
\texttt{ANONYMIZED\_URL} under an open license.
